
In dimensions $n\ge 3$, the corresponding sharp exponent is expected to be
$1/2$ in place of the peak-based exponent $1/(n+2)$ in Theorem~\ref{stabb}.
That higher-dimensional problem remains open and would require a mechanism
beyond the one-dimensional Green-potential and planar free-boundary structure
used above.

%=========================================================
% End of sharpness considerations
%=========================================================

\begin{proof}[Proof of Theorem~\ref{stabb} for $n\ge 3$]
From this point on we use the hyperbolic measure
\[
d\tau(x):=\frac{dV(x)}{(1-|x|^2)^n}.
\]
Fix $n\ge 3$, $\beta>1$, and a measurable set $\Omega\subset\mathbb B^n$ with
\[
0<s_0:=\tau(\Omega)<\infty.
\]
By homogeneity we may assume $\|g\|_{2,\beta}=1$.

Set
\[
u(x):=\Phi_n(x)^\beta g(x)^2,
\qquad
v(x):=\Phi_n(x)^\beta,
\]
and let $u^*,v^*$ be the decreasing rearrangements of $u,v$ with respect to $\tau$.

Since
\[
\|g\|_{2,\beta}^2
=
c_\beta\int_{\mathbb B^n}\Phi_n(x)^\beta g(x)^2\,d\tau(x)=1
\]
and, by the normalization of $c_\beta$,
\[
c_\beta\int_{\mathbb B^n}\Phi_n(x)^\beta\,d\tau(x)=1,
\]
we have
\[
\int_0^\infty u^*(s)\,ds=\int_0^\infty v^*(s)\,ds=\frac1{c_\beta}.
\]

Define
\[
\Theta(s):=\int_0^s v^*(t)\,dt,
\qquad
\theta(s):=c_\beta\,\Theta(s),
\qquad
C(s):=(1-\theta(s))^{-1}.
\]
Equivalently,
\[
\theta(s)
=
c_\beta\int_{B(0,r(s))}\Phi_n(x)^\beta\,d\tau(x),
\]
where $r(s)\in(0,1)$ is determined by
\[
\tau(B(0,r(s)))=s.
\]

\smallskip\noindent
\textbf{1) Deficit on the level set.}
Let
\[
\Omega_{s_0}:=\{u>u^*(s_0)\},
\qquad
\tau(\Omega_{s_0})=s_0.
\]
By Hardy--Littlewood,
\[
\int_\Omega u\,d\tau\le \int_{\Omega_{s_0}}u\,d\tau=\int_0^{s_0}u^*(s)\,ds.
\]
Define
\[
\delta_{s_0}
:=
1-\frac{\displaystyle\int_0^{s_0}u^*(s)\,ds}
{\displaystyle\int_0^{s_0}v^*(s)\,ds}.
\]
Then
\[
\delta_{s_0}\le \delta(g;\Omega).
\]
Also,
\begin{equation}\label{eq:D-def-ng3}
D:=\int_0^{s_0}\bigl(v^*(s)-u^*(s)\bigr)\,ds
=
\delta_{s_0}\int_0^{s_0}v^*(s)\,ds
\le \frac{\delta_{s_0}}{c_\beta}.
\end{equation}

\smallskip\noindent
\textbf{2) The crossing point and the truncated deficit.}
Let
\[
R(s):=\frac{u^*(s)}{v^*(s)}.
\]
For $n\ge3$ we have
\[
(u^*)'(s)+\gamma\,\Upsilon(s)\,u^*(s)\ge0,
\qquad
(v^*)'(s)+\gamma\,\Upsilon(s)\,v^*(s)=0,
\qquad
\gamma:=\beta(n-1)^2.
\]
Hence
\[
R'(s)
=
\frac{(u^*)'(s)v^*(s)-(v^*)'(s)u^*(s)}{(v^*(s))^2}
=
\frac{(u^*)'(s)+\gamma\,\Upsilon(s)\,u^*(s)}{v^*(s)}
\ge0.
\]
Thus $R$ is nondecreasing. Since $u^*(0)\le v^*(0)$ and
\[
\int_0^\infty u^*(s)\,ds=\int_0^\infty v^*(s)\,ds=\frac1{c_\beta},
\]
either $u^*\equiv v^*$ almost everywhere, or else there is a unique point $s^*$ such that
\[
u^*(s^*)=v^*(s^*).
\]
Define
\[
I_*:=\int_0^{s^*}\bigl(v^*(s)-u^*(s)\bigr)\,ds.
\]

We claim that
\begin{equation}\label{eq:Istar-upper-ng3}
I_*\le \frac1{c_\beta}\,C(s_0)\,\delta_{s_0}.
\end{equation}

Arguing exactly as in the case $n=2$, if $s_0\ge s^*$ then
\[
I_*
\le
D\,
\frac{\int_{s^*}^{\infty}u^*(s)\,ds}{\int_{s_0}^{\infty}u^*(s)\,ds}
\le
D\,
\frac{\frac1{c_\beta}}{\int_{s_0}^{\infty}v^*(s)\,ds}.
\]
Since
\[
\int_{s_0}^{\infty}v^*(s)\,ds
=
\frac1{c_\beta}-\Theta(s_0)
=
\frac{1-\theta(s_0)}{c_\beta},
\]
it follows that
\[
I_*
\le
D\,C(s_0)
\le
\frac1{c_\beta}C(s_0)\,\delta_{s_0},
\]
where in the last step we used \eqref{eq:D-def-ng3}.

If $s_0\le s^*$, then, again exactly as in the case $n=2$,
\[
I_*
\le
D\,
\frac{\int_0^{s^*}v^*(s)\,ds}{\int_0^{s_0}v^*(s)\,ds}
\le
D\,
\frac{\frac1{c_\beta}}{\int_0^{s_0}v^*(s)\,ds}
=
\frac{\delta_{s_0}}{c_\beta}
\le
\frac1{c_\beta}C(s_0)\,\delta_{s_0}.
\]
Thus \eqref{eq:Istar-upper-ng3} holds in all cases. Since $\delta_{s_0}\le\delta(g;\Omega)$,
\begin{equation}\label{eq:Istar-upper-final-ng3}
I_*\le \frac1{c_\beta}C(s_0)\,\delta(g;\Omega).
\end{equation}

\smallskip\noindent
\textbf{3) Comparison with the clipped functional $I(T)$.}
Let
\[
T:=\|u\|_{L^\infty(\mathbb B^n)}
=
\sup_{x\in\mathbb B^n}\Phi_n(x)^\beta g(x)^2\in(0,1].
\]
Let $s(T)$ be the unique solution of $v^*(s(T))=T$ and define
\[
I(T):=\int_0^{s(T)}\bigl(v^*(s)-T\bigr)\,ds
=
\Theta(s(T))-T\,s(T).
\]

Since $u^*(s)\le T$ for all $s$ and
\[
u^*(s^*)=v^*(s^*)\le T,
\]
the monotonicity of $v^*$ implies
\[
s^*\ge s(T).
\]
Therefore
\begin{equation}\label{eq:I-le-Istar-ng3}
I(T)
=
\int_0^{s(T)}(v^*(s)-T)\,ds
\le
\int_0^{s(T)}(v^*(s)-u^*(s))\,ds
\le
I_*.
\end{equation}
Combining \eqref{eq:Istar-upper-final-ng3} and \eqref{eq:I-le-Istar-ng3}, we get
\begin{equation}\label{eq:I-upper-deficit-ng3}
I(T)\le \frac1{c_\beta}C(s_0)\,\delta(g;\Omega).
\end{equation}

\smallskip\noindent
\textbf{4) Lower bound on $I(T)$.}
Using $\Theta'(s)=v^*(s)$ and
\[
\frac{\Theta''(s)}{\Theta'(s)}=-\gamma\,\Upsilon(s),
\qquad \gamma=\beta(n-1)^2,
\]
differentiating $v^*(s(T))=\Theta'(s(T))=T$ yields
\[
s'(T)=\frac{1}{\Theta''(s(T))},
\qquad
I'(T)=-s(T),
\qquad
I''(T)=-s'(T)=\frac{1}{\gamma\,T\,\Upsilon(s(T))}.
\]

Fix
\[
r_0^2=\frac1{4(n-2)},\qquad T_0=(1-r_0^2)^{\beta(n-1)}.
\]
For $T\in(T_0,1)$, Lemma~\ref{lem:sT-lower} and Corollary~\ref{cor:Upsilon-sT} give
\[
\Upsilon(s(T))\le
\frac{2\Gamma(\frac n2)\,G_{n,r_0}}{n\pi^{n/2}}
\left(\frac{\beta\,C_{n,r_0}}{1-T}\right)^{\frac{n-2}{2}}.
\]
Hence there exists an explicit constant $\kappa_{n,r_0}>0$ such that
\[
I''(T)\ge \kappa_{n,r_0}\,\beta^{-n/2}\,(1-T)^{\frac{n-2}{2}},
\qquad T\in(T_0,1).
\]
Since $I(1)=I'(1)=0$, integration gives
\[
I(T)\ge
\frac{4\,\kappa_{n,r_0}}{n(n+2)}\,\beta^{-n/2}\,(1-T)^{\frac n2+1},
\qquad T\in(T_0,1).
\]
Set
\[
E_{n,r_0}:=\frac{4\,\kappa_{n,r_0}}{n(n+2)}.
\]
Then
\begin{equation}\label{eq:I-lower-n-ng3}
I(T)\ge E_{n,r_0}\,\beta^{-n/2}\,(1-T)^{\frac n2+1},
\qquad T\in(T_0,1).
\end{equation}

Moreover, $I$ is decreasing on $(0,1)$ because $I'(T)=-s(T)<0$. Hence if $T\le T_0$,
\[
I(T)\ge I(T_0)\ge E_{n,r_0}\,\beta^{-n/2}\,(1-T_0)^{\frac n2+1}.
\]
Since $\beta>1$,
\[
T_0=(1-r_0^2)^{\beta(n-1)}\le (1-r_0^2)^{n-1},
\]
so with
\[
d_n:=1-(1-r_0^2)^{n-1}>0
\]
we get
\begin{equation}\label{eq:I-lower-smallT-ng3}
I(T)\ge E_{n,r_0}\,\beta^{-n/2}\,d_n^{\frac n2+1},
\qquad 0<T\le T_0.
\end{equation}

\smallskip\noindent
\textbf{5) From the deficit to distance from extremals.}
We split into two regimes.

If $T\in(T_0,1)$, then \eqref{eq:I-upper-deficit-ng3} and \eqref{eq:I-lower-n-ng3} imply
\[
E_{n,r_0}\,\beta^{-n/2}\,(1-T)^{\frac n2+1}
\le
\frac1{c_\beta}C(s_0)\,\delta(g;\Omega),
\]
hence
\[
1-T
\le
\left(
\frac{\beta^{n/2}}{c_\beta E_{n,r_0}}\,
C(s_0)\,\delta(g;\Omega)
\right)^{\frac{2}{n+2}}.
\]
By Lemma~\ref{diel},
\[
\inf_{a\in\mathbb B^n}\|g-g_a\|_{2,\beta}^2\le 2-2\sqrt{T}\le 2(1-T),
\]
and therefore
\[
\inf_{a\in\mathbb B^n}\|g-g_a\|_{2,\beta}
\le
\sqrt2
\left(
\frac{\beta^{n/2}}{c_\beta E_{n,r_0}}\,
C(s_0)\,\delta(g;\Omega)
\right)^{\frac1{n+2}}.
\]

If $0<T\le T_0$, then by \eqref{eq:I-upper-deficit-ng3} and \eqref{eq:I-lower-smallT-ng3},
\[
\frac1{c_\beta}C(s_0)\,\delta(g;\Omega)
\ge
E_{n,r_0}\,\beta^{-n/2}\,d_n^{\frac n2+1},
\]
that is,
\[
C(s_0)\,\delta(g;\Omega)
\ge
c_\beta E_{n,r_0}\,\beta^{-n/2}\,d_n^{\frac n2+1}.
\]
Since the left-hand side of the theorem is always bounded by $2$, we obtain
\[
\inf_{\substack{|c|=\|g\|_{2,\beta}\\ a\in\mathbb B^n}}
\frac{\|g-cg_a\|_{2,\beta}}{\|g\|_{2,\beta}}
\le
2\,(c_\beta E_{n,r_0})^{-\frac1{n+2}}\,
\beta^{\frac{n}{2(n+2)}}\,d_n^{-1/2}\,
\bigl(C(s_0)\,\delta(g;\Omega)\bigr)^{\frac1{n+2}}.
\]

Combining the two regimes, the theorem follows with
\[
C_{n,\beta}
=
2(c_\beta E_{n,r_0})^{-\frac1{n+2}}\,
\beta^{\frac{n}{2(n+2)}}
\,d_n^{-1/2},
\qquad
d_n:=1-(1-r_0^2)^{n-1}.
\]
\end{proof}

\begin{numrem}[Explicit form and dimensional dependence of the stability constant]
\label{rem:explicit_Cnalpha}
In the proof of Theorem~\ref{stabb} for $n\ge3$, one may choose
\[
r_0^2=\frac1{4(n-2)},\qquad
G_{n,r_0}:={}_2F_1\!\left(\frac n2,n;\frac{n+2}{2};\,r_0^2\right),
\]
and
\[
C_{n,r_0}:=(n-1)+\frac{(n-1)(n-2)}{n}\,
{}_3F_2\!\left(\begin{array}{c}
1,1,2-\frac n2\\[2pt]
2,1+\frac n2
\end{array};\,r_0^2\right).
\]
Then
\[
\kappa_{n,r_0}
=
\frac{n\,\pi^{n/2}}
{2\,(n-1)^2\,\Gamma(\frac n2)\,G_{n,r_0}\,C_{n,r_0}^{\frac{n-2}{2}}},
\qquad
E_{n,r_0}=\frac{4\,\kappa_{n,r_0}}{n(n+2)}.
\]
With
\[
d_n:=1-(1-r_0^2)^{n-1}>0,
\]
the proof yields
\[
C_{n,\beta}
=
(c_\beta E_{n,r_0})^{-\frac1{n+2}}\,
\beta^{\frac{n}{2(n+2)}}
\max\Bigl\{\sqrt2,\;2d_n^{-1/2}\Bigr\}
=
(c_\beta E_{n,r_0})^{-\frac1{n+2}}\,
\beta^{\frac{n}{2(n+2)}}
\,2d_n^{-1/2}.
\]
In particular, one may write
\[
C_{n,\beta}
=
\widetilde C_n\,
c_\beta^{-\frac1{n+2}}\,
\beta^{\frac{n}{2(n+2)}},
\]
where
\[
\widetilde C_n
:=
2\,E_{n,r_0}^{-\frac1{n+2}}\,d_n^{-1/2}
\]
depends only on $n$.
\end{numrem}

\begin{numrem}[The planar case and the disappearance of the factor $C_{n,r_0}$]
It is worth noting that the higher-dimensional explicit formula is not meant to be specialized directly to $n=2$, since the choice
\[
r_0^2=\frac{1}{4(n-2)}
\]
breaks down at $n=2$. Nevertheless, at a purely formal level one feature is consistent with the planar theory: the exponent
\[
C_{n,r_0}^{\frac{n-2}{2}}
\]
would become
\[
C_{2,r_0}^{0}=1,
\]
so this factor disappears in dimension $2$.

Thus the blow-up of the planar stability constant as the parameter approaches the endpoint does \emph{not} come from the factor $C_{n,r_0}$, but from the normalization of the weighted Bergman norm. More precisely, in the planar notation with parameter $\alpha>-1$ one has
\[
c_\alpha=\frac{\alpha+1}{\pi},
\]
and the stability constant obtained from the separate two-dimensional proof is
\[
C_{2,\alpha}
=
\sqrt2\left(\frac{2(\alpha+2)}{\alpha+1}\right)^{1/4}.
\]
Equivalently, writing $\alpha=\beta-2$, this becomes
\[
C_{2,\beta}
=
\sqrt2\left(\frac{2\beta}{\beta-1}\right)^{1/4},
\]
which coincides with the planar bound obtained earlier.
\end{numrem}
\begin{numlemma}\label{diel}
Let $0<T<1$ and $g$ is log-$\mathcal{M}$-subharmonic function on $\|\cdot\|_{2,\beta}-$norm equal to $1$ such that
$\max_{x \in \mathbb{B}^n}g(x)^2\Phi_n(x)^{{\beta}}=T$. Then,
 $$\inf_{|a|<1}\|g-g_a\|^2\le 2-2 \sqrt{T}.$$
\end{numlemma}
\begin{proof}[Proof of Lemma~\ref{diel}]
Let $m=m_a(x)$ be an involutive Mobius transform of the unit ball onto itself so that $m(0)=a$. Let $$g_a(x)=\frac{\Phi^{\beta/2}_n(m(x))}{{\Phi^{\beta/2}_n(x)}}. $$

Then $$\|g_a\|=1.$$ Let $g$ be log-M-subharmonic so that $\|g\|=1$. Then

\[\begin{split}\|g-g_a\|^2 &= 2 -2\int_{\mathbb{B}} g(x) g_a(x) \Phi_n^\beta(x) d\tau(x)
\\&=2 -2\int_{\mathbb{B}} g(m_a(y)) g_a(m_a(y)) \Phi_n^\beta(m_a(y)) d\tau(m_a(y))
\\&=2 -2\int_{\mathbb{B}} g(m_a(y)) \frac{\Phi^{\beta/2}_n(y)}{{\Phi^{\beta/2}_n(m_a(y))}} \Phi_n^\beta(m_a(y)) d\tau(y)
\\&=2 -2\int_{\mathbb{B}} g(m_a(y)) {\Phi^{\beta/2}_n(y)}\Phi_n^{\beta/2}(m_a(y)) d\tau(y)
\\&\le 2-2 g(a) \Phi_n^{\beta/2}(a).\end{split}\]
In the last inequality we use invariant-mean property of log-$M$-subharmonic function $$h(y)=g(m_a(y)) {\Phi^{-\beta/2}_n(y)}\Phi_n^{\beta/2}(m_a(y))$$ at $y=0$. This follows from \cite[Corollary~4.13]{stoll} by using polar coordinates.

So $$\inf_{|a|<1}\|g-g_a\|^2\le 2-2 \sqrt{T},$$ as desired.\\
\end{proof}

\begin{numlemma}[Explicit lower bound for $s_T$ and a threshold $T_0$]\label{lem:sT-lower}
Assume $n\ge 3$ and $\beta>1$. Set
\[
r_0^2:=\frac{1}{4(n-2)},\qquad T_0:=(1-r_0^2)^{\beta(n-1)}.
\]
For $T\in(T_0,1)$ let $S:=T^{1/\beta}$ and let $r_T\in(0,r_0)$ be uniquely determined by
\[
\Phi_n(r_T)=S .
\]
Then
\[
s_T:=\tau\{x\in\mathbb{B}^n:\ T<\Phi_n^{\beta}(x)<1\}=\tau(\mathbb{B}(0,r_T)).
\]
Moreover, if we denote
\[
H_{n,r_0}:=
{}_3F_2\!\left(\begin{array}{c}
1,1,2-\frac n2\\[2pt]
2,1+\frac n2
\end{array};\, r_0^2\right),
\qquad
C_{n,r_0}:=(n-1)+\frac{(n-1)(n-2)}{n}\,H_{n,r_0},
\]
then for every $T\in(T_0,1)$ we have the quantitative bound
\[
s_T \ge \frac{\omega_n}{n}\left(\frac{1-T}{\beta\,C_{n,r_0}}\right)^{\!n/2},
\]
where $\omega_n:=|\mathbb S^{n-1}|$.
\end{numlemma}

\begin{proof}
Assume $n\ge 3$ and $\beta>1$, and set
\[
r_0^2:=\frac{1}{4(n-2)},\qquad T_0:=(1-r_0^2)^{\beta(n-1)}.
\]
Fix $T\in(T_0,1)$, let $S:=T^{1/\beta}$, and let $r_T\in(0,r_0)$ be determined by
$\Phi_n(r_T)=S$.

Since $\Phi_n$ is radial and strictly decreasing in $r=|x|$ with $\Phi_n(0)=1$,
the superlevel set $\{x\in\mathbb B^n:\ \Phi_n(x)>S\}$ is exactly $\mathbb B(0,r_T)$.
Because $\Phi_n^\beta(x)=\Phi_n(|x|)^\beta$, the condition
$T<\Phi_n^\beta(x)<1$ is equivalent to $\Phi_n(|x|)>S$, hence
\[
\{x\in\mathbb B^n:\ T<\Phi_n^\beta(x)<1\}=\mathbb B(0,r_T),
\]
and therefore
\[
s_T=\tau\{x\in\mathbb B^n:\ T<\Phi_n^{\beta}(x)<1\}=\tau(\mathbb B(0,r_T)).
\]

For $0\le r\le r_0$ we use the derivative identity
\[
-\Phi_n'(r)=2r\left((n-1)+\frac{(n-1)(n-2)}{n}\,
{}_3F_2\!\left(\begin{array}{c}1,1,2-\frac n2\\ 2,1+\frac n2\end{array};\,r^2\right)\right).
\]
The ${}_3F_2$ series has nonnegative coefficients, hence it is increasing in $r^2$,
so for $0\le r\le r_0$,
\[
-\Phi_n'(r)\le 2C_{n,r_0}\,r.
\]
Integrating from $0$ to $r_T$ and using $\Phi_n(0)=1$ yields
\[
1-\Phi_n(r_T)=\int_0^{r_T}(-\Phi_n'(t))\,dt\le C_{n,r_0}r_T^2.
\]
Since $0<\Phi_n(r_T)\le 1$ and $\beta>1$, we have
\[
1-\Phi_n(r_T)^\beta \le \beta\,(1-\Phi_n(r_T)),
\]
and because $\Phi_n(r_T)^\beta=T$ this gives
\[
1-T \le \beta\,C_{n,r_0}\,r_T^2,
\qquad\text{hence}\qquad
r_T^2\ge \frac{1-T}{\beta\,C_{n,r_0}}.
\]

Finally, by definition of $\tau$,
\[
\tau(\mathbb B(0,r_T))
=
\omega_n\int_0^{r_T}\frac{\rho^{n-1}}{(1-\rho^2)^n}\,d\rho
\ge
\omega_n\int_0^{r_T}\rho^{n-1}\,d\rho
=
\frac{\omega_n}{n}\,r_T^n,
\]
so using the lower bound for $r_T^2$ we obtain
\[
s_T=\tau(\mathbb B(0,r_T))
\ge \frac{\omega_n}{n}\left(\frac{1-T}{\beta\,C_{n,r_0}}\right)^{n/2}.
\]
This proves the lemma.
\end{proof}

\begin{numcor}[Explicit bound for $\Upsilon(s_T)$]\label{cor:Upsilon-sT}
With the same assumptions and notation as in Lemma~\ref{lem:sT-lower}, set
\[
G_{n,r_0}:={}_2F_1\!\left(\frac n2,n;\frac{n+2}{2};\,r_0^2\right).
\]
Then for all $T\in(T_0,1)$,
\[
\Upsilon(s_T)\le
\frac{2\Gamma\!\left(\frac n2\right) G_{n,r_0}}{n\pi^{n/2}}
\left(\frac{\beta\,C_{n,r_0}}{1-T}\right)^{\!\frac{n-2}{2}}.
\]
\end{numcor}

\begin{proof}
With the notation of Lemma~\ref{lem:sT-lower}, we have $s_T=\tau(\mathbb B(0,r_T))$.
By the definition of $\Upsilon$ in terms of the radius $r=r(s)$ determined by
$s=\tau(\mathbb B(0,r))$,
\[
\Upsilon(s)=\frac{2\Gamma\!\left(\frac n2\right)}{n\pi^{n/2}}\,
r(s)^{\,2-n}\,
{}_2F_1\!\left(\frac n2,n;\frac{n+2}{2};\,r(s)^2\right),
\]
hence
\[
\Upsilon(s_T)=\frac{2\Gamma\!\left(\frac n2\right)}{n\pi^{n/2}}\,
r_T^{\,2-n}\,
{}_2F_1\!\left(\frac n2,n;\frac{n+2}{2};\,r_T^2\right).
\]
The ${}_2F_1$ series has positive coefficients, so it is increasing in $r^2$; since
$r_T\le r_0$ we obtain
\[
{}_2F_1\!\left(\frac n2,n;\frac{n+2}{2};\,r_T^2\right)
\le {}_2F_1\!\left(\frac n2,n;\frac{n+2}{2};\,r_0^2\right)
=G_{n,r_0}.
\]
Therefore
\[
\Upsilon(s_T)\le
\frac{2\Gamma\!\left(\frac n2\right) G_{n,r_0}}{n\pi^{n/2}}\,
r_T^{\,2-n}.
\]
Finally, Lemma~\ref{lem:sT-lower} gives
\[
r_T^2\ge \frac{1-T}{\beta\,C_{n,r_0}},
\]
hence
\[
r_T^{\,2-n}\le \left(\frac{\beta\,C_{n,r_0}}{1-T}\right)^{\frac{n-2}{2}},
\]
which yields
\[
\Upsilon(s_T)\le
\frac{2\Gamma\!\left(\frac n2\right) G_{n,r_0}}{n\pi^{n/2}}
\left(\frac{\beta\,C_{n,r_0}}{1-T}\right)^{\!\frac{n-2}{2}}.
\]
This proves the corollary.
\end{proof}




%----------------------------------------------------------
% Remark (optional): you can add a short paragraph here saying that
% families satisfying \eqref{eq:clipped-assumption} are obtained by
% ``max-type'' or ``log-sum-exp'' combinations of two far-apart extremals,
% e.g. q->infty of (g_{a_1}^q + \lambda^q g_{a_2}^q)^{1/q}, with separation.
%----------------------------------------------------------








\section{Log-subharmonicity and Fock space}

\begin{numdefn}
Let $f$ be a real-analytic complex function defined in $\mathbb{R}^m$, $m\ge 2$, whose absolute value $|f|$ is log-harmonic in $\mathbb{R}^m\setminus f^{-1}(0)$.
We say that $f$ Belongs to the Fock space $\Fock$ if $$\|f\|=\left(\int_{\mathbb{R}^n}|f(x)|^2 e^{-\pi |x|^2} dx\right)^{1/2}<\infty.$$
\end{numdefn}

Let $g_{z_0}=|f_{z_0}|$, where $f_{z_0}=e^{-\frac{\pi}{2}|z_0|^2}e^{\pi z\overline{z_0}}.$
\begin{numthm}[Fock space version of Theorem \ref{stabb}]\label{thm:StabilityFockSpace}
    There is an explicitly computable constant $C>0$ such that, for all  measurable sets $\Om \subset \mathbb{C}$ with finite measure and for the set $\Fock$ of all log-subharmonic  real-analytic functions in $\C$ belonging to the Fock space and $F\in \Fock$, we have
    \begin{equation}\label{eqn:StabilityFunctionFock}
        \min_{\substack{c=\|F\|_{\Fock},\\ z_0\in \C}} \frac{\Fnorm{F-cg_{z_0}}}{\Fnorm{F}} \leq C
        \left(e^{\vol{\Om}} \delta(F;\Om)\right)^{1/4},
    \end{equation}
where
\begin{equation}
\label{defdeltaFO}
\delta(F;\Omega) \coloneqq 1-\frac{\int_\Omega |F(z)|^2 e^{-\pi |z|^2}\diff z}
{(1-e^{-|\Omega|})\Vert F\Vert_{\Fock}^2}.
\end{equation}

\end{numthm}
\begin{numlemma}\label{ulqin}
Assume that $f:\mathbb C\to[0,\infty)$ is a real-analytic log-subharmonic function
belonging to the Fock space, i.e.
\[
\|f\|_{\Fock}^2:=\int_{\mathbb C} f(w)^2 e^{-\pi |w|^2}\,dA(w)<\infty.
\]
Then for every $z\in\mathbb C$,
\[
f(z)^2e^{-\pi |z|^2}\le \|f\|_{\Fock}^2.
\]
Moreover, if $\|f\|_{\Fock}=1$, then
\[
\inf_{z_0\in \mathbb C}\|f-g_{z_0}\|_{\Fock}^2\le 2-2 \sqrt{T},
\]
where
\[
T=\sup_{z\in \mathbb C} f(z)^2e^{-\pi |z|^2},
\qquad
g_{z_0}(z)=e^{-\frac{\pi}{2}|z_0|^2+\pi \Re(z\overline{z_0})}.
\]
\end{numlemma}

\begin{proof}
Fix \(z\in\mathbb C\), and define
\[
H_z(w):=f(z+w)^2 e^{-2\pi\Re((w+z)\overline z)}.
\]
Since \(f\) is positive and log-subharmonic, \(f^2\) is also log-subharmonic, hence
\(H_z\) is log-subharmonic as well. Applying the Gaussian mean-value inequality to \(H_z\)
at the origin, we obtain
\[
H_z(0)\le \int_{\mathbb C} H_z(w)e^{-\pi|w|^2}\,dA(w).
\]
Now
\[
H_z(0)=f(z)^2e^{-2\pi|z|^2},
\]
while
\[
H_z(w)e^{-\pi|w|^2}
=
f(z+w)^2 e^{-2\pi\Re((w+z)\overline z)}e^{-\pi|w|^2}.
\]
Changing variables \(u=z+w\), we get
\[
f(z)^2e^{-2\pi|z|^2}
\le
\int_{\mathbb C}f(u)^2 e^{-2\pi\Re(u\overline z)}e^{-\pi|u-z|^2}\,dA(u).
\]
Using
\[
-2\Re(u\overline z)-|u-z|^2=-|u|^2+|z|^2,
\]
it follows that
\[
f(z)^2e^{-2\pi|z|^2}
\le
e^{-\pi|z|^2}\int_{\mathbb C}f(u)^2e^{-\pi|u|^2}\,dA(u),
\]
and hence
\[
f(z)^2e^{-\pi|z|^2}\le \|f\|_{\Fock}^2.
\]

Assume now that \(\|f\|_{\Fock}=1\). Since also \(\|g_{z_0}\|_{\Fock}=1\), we have
\[
\|f-g_{z_0}\|_{\Fock}^2
=
2-2\langle f,g_{z_0}\rangle_{\Fock},
\]
where
\[
\langle f,g_{z_0}\rangle_{\Fock}
=
\int_{\mathbb C} f(z)g_{z_0}(z)e^{-\pi|z|^2}\,dA(z).
\]
Now the function \(h_{z_0}:=f\,g_{z_0}\) is again positive and log-subharmonic, since
\(g_{z_0}\) is positive and log-harmonic. Applying the first part of the lemma to \(h_{z_0}\),
we obtain
\[
f(z_0)g_{z_0}(z_0)e^{-\pi|z_0|^2}
\le
\int_{\mathbb C} f(z)g_{z_0}(z)e^{-\pi|z|^2}\,dA(z)
=
\langle f,g_{z_0}\rangle_{\Fock}.
\]
Since
\[
g_{z_0}(z_0)=e^{-\frac{\pi}{2}|z_0|^2+\pi|z_0|^2}
=e^{\frac{\pi}{2}|z_0|^2},
\]
we get
\[
\langle f,g_{z_0}\rangle_{\Fock}\ge f(z_0)e^{-\frac{\pi}{2}|z_0|^2}.
\]
Therefore
\[
\|f-g_{z_0}\|_{\Fock}^2
\le
2-2\,f(z_0)e^{-\frac{\pi}{2}|z_0|^2}.
\]
Taking the infimum over \(z_0\in\mathbb C\), and using the definition of \(T\), we conclude that
\[
\inf_{z_0\in\mathbb C}\|f-g_{z_0}\|_{\Fock}^2
\le
2-2\sup_{z_0\in\mathbb C} f(z_0)e^{-\frac{\pi}{2}|z_0|^2}
=
2-2\sqrt{T}.
\]
\end{proof}
\begin{proof}[Schetch of the proof of Theorem~\ref{thm:StabilityFockSpace}]
The proof if just an imitation of the corresponding theorem in \cite{GomezGuerraRamosTilli2024}. Instead of holomorphicity we have more general log-subharmonicity of our function $F$.
\begin{equation}
\label{defu}
u(z):=F^2(z) e^{-\pi|z|^2},
\end{equation}
in connection with its super-level sets



\begin{equation}
\label{defAt}
A_t:=\{u>t\}=\left\{z\in\mathbb{C}\,:\,\, u(z)>t\right\},
\end{equation}
its \emph{distribution function}




\begin{equation}
\label{defmu}
\mu(t):= |A_t|,\qquad 0\leq t\leq \max_{\C} u
\end{equation}
(note that $u$ is bounded due to \eqref{ulqin}),
and the \emph{decreasing rearrangement} of $u$, i.e. the function
\begin{equation}
\label{defclassu*}
u^*(s):=\sup\{t\geq 0\,:\,\, \mu(t)>s\}\qquad \text{for $s\geq 0$}
\end{equation}
(for more details on rearrangements, we refer to  \cite{baernstein}).
Then in \cite{Nicola} it is proved that \begin{equation}\label{nic}(u^\ast)'(s)+u^\ast(s)\ge 0. \end{equation}
The same proof gives the same inequality for $F$ being a log-subharmonic function in the complex plane. Now the proof of our theorem follows from Lemma~\ref{ulqin}, \eqref{nic}, by following the same approach as in \cite{GomezGuerraRamosTilli2024}. Namely  Lemma~2.3 from  \cite{GomezGuerraRamosTilli2024} and  eq.~246. from \cite{GomezGuerraRamosTilli2024} follows from \eqref{nic}.

\end{proof}

By using the following lemma, we can state our stability result for the Fock space in higher dimensions.
\begin{numlemma}\label{ulqin2024}
Assume that \(h\) is a real-analytic log-subharmonic function in \(\mathbb{R}^m\) belonging to the Fock-type space. Then for every \(x\in \mathbb{R}^m\),
\begin{equation}\label{sharine1}
h(x)e^{-\pi |x|^2}\le  \int_{\mathbb{R}^m} h(y) e^{-\pi |y|^2}\, dA(y).
\end{equation}
Assume moreover that \(f\) is a real-analytic log-subharmonic function in \(\mathbb{R}^m\) belonging to the Fock space, that \(\|f\|_{\Fock}=1\), and define
\[
g_{x_0}(y)= e^{-\frac{\pi}{2} |x_0|^2+\pi \langle y,x_0\rangle}.
\]
Then
\begin{equation}\label{ulqin1}
\inf_{x_0\in \mathbb{R}^m}\|f-g_{x_0}\|^2\le 2-2 \sqrt{T},
\qquad
T=\sup_{x\in \mathbb{R}^m} |f(x)|^2e^{-\pi |x|^2}.
\end{equation}
\end{numlemma}

\begin{proof}
Fix \(x\in \mathbb{R}^m\), and consider the function
\[
H_x(y):=h(x+y)e^{-2\pi \langle x+y,x\rangle}, \qquad y\in \mathbb{R}^m.
\]
Since \(h\) is log-subharmonic, \(\log h\) is subharmonic wherever \(h>0\). The function
\[
y\mapsto -2\pi \langle x+y,x\rangle
\]
is affine, hence harmonic, so \(\log H_x=\log h(x+y)-2\pi\langle x+y,x\rangle\) is again subharmonic. Therefore \(H_x\) is log-subharmonic, and in particular subharmonic.

Applying the Gaussian mean-value inequality to \(H_x\) at the origin gives
\[
H_x(0)\le \int_{\mathbb{R}^m} H_x(y)e^{-\pi |y|^2}\,dA(y).
\]
Because
\[
H_x(0)=h(x)e^{-2\pi |x|^2},
\]
we obtain
\[
h(x)e^{-2\pi |x|^2}
\le
\int_{\mathbb{R}^m} h(x+y)e^{-2\pi\langle x+y,x\rangle}e^{-\pi|y|^2}\,dA(y).
\]
Now make the change of variables \(u=x+y\). Then \(y=u-x\), and the right-hand side becomes
\[
\int_{\mathbb{R}^m} h(u)e^{-2\pi\langle u,x\rangle}e^{-\pi|u-x|^2}\,dA(u).
\]
Using the identity
\[
2\langle u,x\rangle+|u-x|^2=|u|^2+|x|^2,
\]
we get
\[
e^{-2\pi\langle u,x\rangle}e^{-\pi|u-x|^2}
=
e^{-\pi|u|^2}e^{-\pi|x|^2}.
\]
Hence
\[
h(x)e^{-2\pi |x|^2}
\le
e^{-\pi|x|^2}\int_{\mathbb{R}^m} h(u)e^{-\pi|u|^2}\,dA(u),
\]
and multiplying by \(e^{\pi|x|^2}\) yields
\[
h(x)e^{-\pi |x|^2}\le  \int_{\mathbb{R}^m} h(u)e^{-\pi |u|^2}\,dA(u).
\]
This proves \eqref{sharine1}.

Now let \(x_0\in \mathbb{R}^m\). Since \(\|f\|_{\Fock}=1\) and \(\|g_{x_0}\|_{\Fock}=1\), we have
\[
\|f-g_{x_0}\|^2
=
\|f\|_{\Fock}^2+\|g_{x_0}\|_{\Fock}^2
-2\int_{\mathbb{R}^m} f(y)g_{x_0}(y)e^{-\pi|y|^2}\,dA(y),
\]
that is,
\[
\|f-g_{x_0}\|^2
=
2-2\int_{\mathbb{R}^m} f(y)g_{x_0}(y)e^{-\pi|y|^2}\,dA(y).
\]

The function \(y\mapsto f(y)g_{x_0}(y)\) is again log-subharmonic: indeed, \(f\) is log-subharmonic and
\[
\log g_{x_0}(y)= -\frac{\pi}{2}|x_0|^2+\pi\langle y,x_0\rangle
\]
is affine, hence harmonic. Therefore we may apply \eqref{sharine1} with \(h(y)=f(y)g_{x_0}(y)\). This gives
\[
f(x_0)g_{x_0}(x_0)e^{-\pi|x_0|^2}
\le
\int_{\mathbb{R}^m} f(y)g_{x_0}(y)e^{-\pi|y|^2}\,dA(y).
\]
Since
\[
g_{x_0}(x_0)
=
e^{-\frac{\pi}{2}|x_0|^2+\pi|x_0|^2}
=
e^{\frac{\pi}{2}|x_0|^2},
\]
the left-hand side simplifies to
\[
f(x_0)e^{-\frac{\pi}{2}|x_0|^2}.
\]
Therefore
\[
f(x_0)e^{-\frac{\pi}{2}|x_0|^2}
\le
\int_{\mathbb{R}^m} f(y)g_{x_0}(y)e^{-\pi|y|^2}\,dA(y),
\]
and substituting this into the norm identity gives
\[
\|f-g_{x_0}\|^2
\le
2-2f(x_0)e^{-\frac{\pi}{2}|x_0|^2}.
\]
Since this holds for every \(x_0\in\mathbb{R}^m\), we obtain
\[
\inf_{x_0\in\mathbb{R}^m}\|f-g_{x_0}\|^2
\le
2-2\sup_{x\in\mathbb{R}^m} f(x)e^{-\frac{\pi}{2}|x|^2}.
\]
Because a log-subharmonic function is nonnegative, we may write
\[
\sup_{x\in\mathbb{R}^m} f(x)e^{-\frac{\pi}{2}|x|^2}
=
\left(\sup_{x\in\mathbb{R}^m}|f(x)|^2e^{-\pi|x|^2}\right)^{1/2}
=
\sqrt{T}.
\]
Hence
\[
\inf_{x_0\in\mathbb{R}^m}\|f-g_{x_0}\|^2\le 2-2\sqrt{T},
\]
which is exactly \eqref{ulqin1}.
\end{proof}
Let $\mu(t)=|A_t|$, where $A_t=\{x\in\mathbb{R}^m: f^2(x)e^{-\pi |x|^2}>t\}$ and
$u^*\colon \R^+\to (0,T]$ is the decreasing rearrangement of $f$, usually
defined as
\begin{equation}\label{eqn:defu*}
    u^*(s) \coloneqq \sup \inset{t \geq 0 \colon \mu(t) \gs s},\qquad
    s\geq 0.
\end{equation}
Similarly let $v$ be the the decreasing rearrangement of $f$ of $e^{-\pi/2 |x|^2}.$ Then $$v^\ast(s)=e^{-\left(s \Gamma\left[1+\frac{m}{2}\right]\right)^{2/m}}$$
By proving the following differential inequality for the decreasing rearrangement $u^*(s)$ (i.e.,\ the inverse function of $\mu(t)$) in the same way as in \cite{GomezGuerraRamosTilli2024}:
\begin{equation}
    \label{diffinequd}
    (u^*)'(s)+\frac {2(\Gamma[m/2+1]\,!)^{\frac {2}{m}}\,\, u^*(s)}{m \,\, s^{1-\frac {2}{m}}}\geq 0,\quad \text{for a.e. $s>0$.}
\end{equation}
we can deduce the following counterpart of \cite{GomezGuerraRamosTilli2024}[7.24]
\begin{equation}
    \label{estlemma1d}
\frac {(1-T)^{m/2+1}}{\Gamma((m/2+2))}\leq \int_0^{s^*}\left(\ee(s)-u^*(s)\right)\diff s\leq
\frac {\delta_{s_0}}{\int_{s_0}^\infty \ee(s)\diff s}.
\end{equation} Then in view of Lemma~\ref{ulqin2024}, in the same way as in \cite{GomezGuerraRamosTilli2024} we can prove the following theorem
\begin{numthm}[Higher dimensional version of Theorem \ref{thm:StabilityFockSpace}]\label{thm:StabilityFockSpaceDimensiond}
   There is a computable constant $C=C(m)>0$ such that, for all  measurable sets $\Om \subset \R^{m}$ with finite measure $|\Omega|>0$ and all real-analytic log-subharmonic functions $f\in \Fock(\R^m)\backslash\{0\}$, we have
    \begin{equation}\label{eqn:StabilityFunctionFockDimensiond}
        \min_{\substack{|c|=\|f\|_{\Fock},\\ x_0\in \R^m}} \frac{\Fnorm{f-cg_{x_0}}}{\Fnorm{f}} \leq
        C
        \left(\frac{\delta(f;\Om)}
      {\int_{|\Omega|}^\infty   e^{-\left(s \Gamma\left[1+\frac{m}{2}\right]\right)^{2/m}}\diff s  } \right)^{\frac 1 {2+m}},
    \end{equation}
where
\begin{equation}
\label{defdeltaFOd}
\delta(f;\Omega):=1-\frac{\int_\Omega |f(x)|^2 e^{-\pi |x|^2}\diff x}
{\Vert f\Vert_{\Fock}^2\, \int_0^{|\Omega|} \ee(s)\diff s}.
\end{equation}
\end{numthm}

\subsection*{Acknowledgements} The second author has been partially supported by the Ministry of Education, Science and Innovation of Montenegro through the grant ``Mathematical Analysis, Optimization and Machine Learning'' and "Complex-analytic and geometric techniques for non-Euclidean machine learning: theory and applications". The third author is partially supported by MPNTR grant 174017.

%  Self-contained Hilbert-envelope / modulus stability corollary
%  (with an explicit vector-valued reduction and citations)
% =========================
% =========================
%  Minimal LaTeX preamble + common definitions
%  (copy/paste as a standalone header)
% =========================




% ==========================================================
%  Theorem environments (edit as you like)
% ==========================================================


% ==========================================================
%  Document starts here
% ==========================================================




\begin{thebibliography}{99}

\bibitem{14}
E.~Schmidt, \emph{{\"U}ber die isoperimetrische {A}ufgabe im $n$-dimensionalen {R}aum konstanter negativer Kr{\"u}mmung. I. Die isoperimetrischen {U}ngleichungen in der hyperbolischen {E}bene und f\"ur {R}otationsk{\"o}rper im $n$-dimensionalen hyperbolischen {R}aum}, \emph{Mathematische Zeitschrift} \textbf{46} (1940), 204--230.

\bibitem{baernstein}
A.~Baernstein, II, \emph{Symmetrization in analysis}, volume~36 of \emph{New Mathematical Monographs}, Cambridge University Press, Cambridge, 2019. With David Drasin and Richard S.~Laugesen, with a foreword by Walter Hayman.

\bibitem{cvpde}
V.~B\"ogelein, F.~Duzaar, and C.~Scheven, \emph{A sharp quantitative isoperimetric inequality in hyperbolic $n$-space}, \emph{Calculus of Variations and Partial Differential Equations} \textbf{54} (2015), no.~5, 3967--4017.

\bibitem{GomezKalajMelentijevicRamos2025}
J.~G{\'o}mez, D.~Kalaj, P.~Melentijevi{\'c}, and J.~P.~G. Ramos, \emph{Uniform stability of concentration inequalities and applications}, \emph{Proc. Lond. Math. Soc. (3)} \textbf{131} (2025), no.~6, Article ID e70114, 53~pp.

\bibitem{GomezGuerraRamosTilli2024}
J.~G{\'o}mez, A.~Guerra, J.~P.~G. Ramos, and P.~Tilli, \emph{Stability of the Faber--Krahn inequality for the short-time Fourier transform}, \emph{Invent. Math.} \textbf{236} (2024), no.~2, 779--836.

\bibitem{kalaj2024}
D.~Kalaj, \emph{Contraction property of certain classes of log-$M$-subharmonic functions in the unit ball}, \emph{J. Funct. Anal.} \textbf{286} (2024), no.~1, Article ID 110203, 29~pp.

\bibitem{kalajramos2024}
D.~Kalaj and J.~P.~G. Ramos, \emph{A Faber--Krahn type inequality for log-subharmonic functions in the hyperbolic ball}, \emph{Israel J. Math.} (2025), published online 29 June 2025. \textsc{doi}:10.1007/s11856-025-2788-0. Preprint: arXiv:2303.08069.

\bibitem{kulikov2022}
A.~Kulikov, \emph{Functionals with extrema at reproducing kernels}, \emph{Geom. Funct. Anal.} \textbf{32} (2022), no.~4, 938--949. \textsc{doi}:10.1007/s00039-022-00608-5.

\bibitem{KulikovNicolaOrtegaCerdaTilli2025}
A.~Kulikov, F.~Nicola, J.~Ortega-Cerd\`a, and P.~Tilli, \emph{A monotonicity theorem for subharmonic functions on manifolds}, \emph{Adv. Math.} \textbf{479} (2025), 110423. \textsc{doi}:10.1016/j.aim.2025.110423. Preprint: arXiv:2212.14008.

\bibitem{LiebLoss}
E.~H. Lieb and M.~Loss, \emph{Analysis}, 2nd ed., Graduate Studies in Mathematics, Vol.~14, American Mathematical Society, Providence, RI, 2001.

\bibitem{Melentijevic}
P. Melentijevi\'c \emph{Sharp stability of convex functionals on weighted Bergman spaces with applications},  arXiv:2508.01939v4


\bibitem{Nicola}
F.~Nicola and P.~Tilli, \emph{The Faber--Krahn inequality for the short-time Fourier transform}, \emph{Inventiones mathematicae} \textbf{230} (2022), 1--30.

\bibitem{RamosTilli}
J.~P.~G. Ramos and P.~Tilli, \emph{A Faber--Krahn inequality for wavelet transforms}, \emph{Bull. Lond. Math. Soc.} \textbf{55} (2023), no.~4, 2018--2034. \textsc{doi}:10.1112/blms.12833.


\bibitem{rudin}
W.~Rudin, \emph{Function theory in the unit ball of $\mathbb{C}^n$}, reprint of the 1980 original, Classics in Mathematics, Springer, Berlin, 2008.

\bibitem{stoll}
M.~Stoll, \emph{Harmonic and subharmonic function theory on the hyperbolic ball}, London Mathematical Society Lecture Note Series~431, Cambridge University Press, Cambridge, 2016.

\end{thebibliography}
\end{document}
%\printbibliography[heading=bibintoc,title=References]        % References and toc inclusion

%\nocite{*}                                  % Include all the resources
                        % References and toc inclusion


% ============================================================
% 4. Möbius-transfer to nonradial data
% ============================================================

\begin{numcor}[Hyperbolic-radial transfer]
\label{cor:mobius_transfer}
Let $\alpha>-1$ and let $g:\D\to[0,\infty)$ be log-subharmonic with
\[
\|g\|_{2,\alpha}=1.
\]
Assume there exists $a\in\D$ such that
\[
\Phi_\alpha(z):=g(z)^2(1-|z|^2)^{\alpha+2}
\]
depends only on $\rho_a(z):=|\varphi_a(z)|$ and is nonincreasing as a function of $\rho_a$,
where
\[
\varphi_a(z)=\frac{a-z}{1-\overline a z}.
\]
Let $\Omega\subset\D$ be measurable with $\mu(\Omega)=s$, and define
\[
g_a(z):=|\varphi_a'(z)|^{(\alpha+2)/2}.
\]
Then
\[
\inf_{b\in\D}\|g-g_b\|_{2,\alpha}
\le \sqrt{3}\,\bigl(C_\alpha(s)\,\delta(g;\Omega,\alpha)\bigr)^{1/2},
\]
where
\[
\delta(g;\Omega,\alpha)
:=1-\frac{1}{\theta_\alpha(s)}\int_\Omega g(z)^2\,dA_\alpha(z).
\]
\end{numcor}

\begin{proof}
Define
\[
(T_ag)(z):=g(\varphi_a(z))\,|\varphi_a'(z)|^{(\alpha+2)/2}.
\]
Using
\[
1-|\varphi_a(z)|^2 = |\varphi_a'(z)|(1-|z|^2),
\qquad
dA(\varphi_a(z))=|\varphi_a'(z)|^2\,dA(z),
\]
a change of variables shows that $T_a$ is an isometry on $L^2(dA_\alpha)$:
\[
\|T_ag\|_{2,\alpha}=\|g\|_{2,\alpha}=1.
\]
Also, $\log(T_ag)=\log g\circ \varphi_a + \frac{\alpha+2}{2}\log|\varphi_a'|$ is subharmonic,
because $\log g$ is subharmonic and $\log|\varphi_a'|$ is harmonic. Hence $T_ag$ is log-subharmonic.

Next,
\[
(T_ag)(z)^2(1-|z|^2)^{\alpha+2}
=
g(\varphi_a(z))^2(1-|\varphi_a(z)|^2)^{\alpha+2}
=
\Phi_\alpha(\varphi_a(z)).
\]
By hypothesis, this depends only on $|z|$ and is nonincreasing in $|z|$.
Therefore Proposition~\ref{prop:bathtub_monotone} applies to $T_ag$.

Since $d\mu$ is Möbius invariant,
\[
\mu(\varphi_a^{-1}(\Omega))=\mu(\Omega)=s,
\]
and another change of variables gives
\[
\int_{\varphi_a^{-1}(\Omega)} (T_ag)^2\,dA_\alpha
=
\int_\Omega g^2\,dA_\alpha.
\]
Hence
\[
\delta(T_ag;\varphi_a^{-1}(\Omega),\alpha)=\delta(g;\Omega,\alpha).
\]
Applying Proposition~\ref{prop:bathtub_monotone} to $T_ag$ and $\varphi_a^{-1}(\Omega)$, we get
\[
\|T_ag-1\|_{2,\alpha}
\le
\sqrt{3}\,\bigl(C_\alpha(s)\,\delta(g;\Omega,\alpha)\bigr)^{1/2}.
\]

Finally, for
\[
g_a(z)=|\varphi_a'(z)|^{(\alpha+2)/2},
\]
one has $T_a g_a\equiv 1$ (because $\varphi_a\circ\varphi_a=\mathrm{id}$ for this automorphism),
so by isometry,
\[
\|g-g_a\|_{2,\alpha}=\|T_ag-1\|_{2,\alpha}.
\]
Therefore
\[
\inf_{b\in\D}\|g-g_b\|_{2,\alpha}
\le \|g-g_a\|_{2,\alpha}
\le \sqrt{3}\,\bigl(C_\alpha(s)\,\delta(g;\Omega,\alpha)\bigr)^{1/2}.
\]
\end{proof}
